\section{Introduction}
\label{sec:introduction}

Remove 5\% of vision tokens from SmolVLA and its success rate collapses from 59.2\% to 0.0\% across all task suites and pruning methods. Apply the same 5\% reduction to $\pi_0.5$ and it retains 98.2\%. Both are Vision-Language-Action (VLA) models built on a Vision Transformer (ViT) encoder feeding tokens through a connector into a language model backbone; the critical difference is not model scale (450M vs.\ 4B parameters) but the connector architecture. Token pruning is well-established for image classification, where DynamicViT~\cite{rao2021dynamicvit}, Token Merging~\cite{bolya2023tome}, and EViT~\cite{liang2022evit} reduce computation by 50\% or more with minimal accuracy loss. A growing body of work extends these techniques to VLAs: VLA-IAP~\cite{wang2026vlaiap}, VLA-Cache~\cite{hsu2025vlacache}, EfficientVLA~\cite{yang2025efficientvla}, and VLA-Pruner~\cite{liu2025vlapruner} each demonstrate efficiency gains on individual architectures.

These methods implicitly assume that token pruning effectiveness transfers across VLA designs. Yet current VLAs employ fundamentally different vision connectors: linear projections~\cite{black2024pi0}, multi-layer MLP projectors~\cite{kim2024openvla}, pixel-shuffle spatial folding~\cite{shukor2025smolvla}, and learned encoders with cross-token self-attention~\cite{zheng2025xvla}. Each connector imposes distinct constraints on how visual information is distributed across tokens. Deploying a pruning method validated on one architecture to a different connector type risks silent degradation or, as the SmolVLA result above illustrates, complete loss of manipulation capability. No prior work systematically evaluates pruning robustness across both connector architectures and task complexity distributions, and the factor that determines whether a given VLA tolerates token removal remains uncharacterized.

We propose a two-axis analysis of VLA token compressibility: pruning sensitivity as a function of connector coupling and task complexity (Figure~\ref{fig:teaser}). Along the coupling axis, token-independent connectors (linear projections, MLP projectors) degrade gracefully because each output token encodes information from a single input patch. Token-entangling connectors produce qualitatively different failures: SmolVLA's pixel-shuffle ($4{\times}4$ spatial folding) collapses to 0.0\% success rate at every retention ratio tested, while X-VLA's 12-layer BART encoder exhibits sharp, threshold-dependent drops. Along the task complexity axis, multi-step tasks amplify degradation in a connector-dependent manner. At retention ratio $k{=}0.9$, X-VLA loses 1.6 percentage points on single-step Object manipulation but 27.8 percentage points on multi-step Goal-conditioned tasks, a gap driven entirely by task demands at the same retention ratio. The coefficient of variation (CV) of post-connector token L2 norms captures connector coupling on a continuous spectrum from 0.038 ($\pi_0$, linear projection) to 0.314 (SmolVLA, pixel-shuffle) and correlates with both pruning tolerance and the retention threshold at which pruning method choice becomes critical. At $k{=}0.95$, $\pi_0.5$ (CV$\,{=}\,$0.055) exhibits a 28.8 percentage point gap between L2 norm selection and random subsampling on LIBERO-Object, while X-VLA and OpenVLA (CV$\,{\geq}\,$0.15) show gaps below 1.5 percentage points, a 24$\times$ difference in method sensitivity.

For pixel-shuffle connectors where all pruning strategies yield 0.0\% success rate, we introduce Vision Token Re-encoder (VTR), a lightweight cross-attention module (${\sim}$22M parameters, approximately 5\% of SmolVLA) that recovers over 60\% success rate through learned re-encoding at 75\% token compression. The VTR performance curve across output token counts $M$ is flat ($M{=}16$ matches $M{=}128$ within 1.1 percentage points), which confirms that learned re-encoding of spatially entangled representations, not the compression ratio, is the operative factor.

Our contributions are:
\begin{itemize}
    \item A two-axis compressibility analysis validated across six VLA models (450M to 7B parameters) on four LIBERO~\cite{liu2023libero} task suites, which reveals connector architecture and task complexity as interacting determinants of pruning robustness.
    \item CV of token L2 norms as a single-pass diagnostic that correlates with both pruning tolerance and method sensitivity, and enables connector-appropriate compression strategy selection without downstream evaluation.
    \item Vision Token Re-encoder (VTR), a learned re-encoding module that recovers pixel-shuffle VLAs from 0.0\% to over 60\% success rate at 75\% token compression, where all hard pruning methods fail.
\end{itemize}
